\subsection{Alignment limits scaling}\label{sec:disc:small}

\paragraph{Alignment limits scaling.} On primates14.chr6 every stage except
alignment keeps getting faster up to 64 ranks, while alignment stops
improving at about 1,200\,s. Table~\ref{tab:perrank} points to how the
work is assigned. The average rank's alignment time halves each time the
ranks double, as it should, but the slowest rank's barely changes, and the
alignment stage lasts as long as that slowest rank. The work is therefore
divisible in total but unevenly assigned: at 64 ranks the fastest rank
finishes in 32\,s and waits about 20 minutes for the slowest. We think the cause is the unit of assignment. A rank aligns all seeds of
the contig pairs it owns, and a pair is never split, so one expensive pair
sets a lower bound on the alignment stage.

\paragraph{Weak scaling.} At a fixed number of bases per rank, seed
generation and redistribution stay flat, and extraction and the sort grow
modestly. Alignment grows
18$\times$ from yeast7 to primates14.chr6, for two reasons
(Table~\ref{tab:weak}).

First, equal input does not mean equal alignment work. The work depends on
how much of the input matches other input: a stretch of one assembly that
matches 13 others is aligned 13 times. The 14 primate assemblies match each
other almost end to end, so each rank aligns 7.9$\times$ as many bases as on
yeast7, and the average rank's alignment time grows by a similar
8.6$\times$. The number of alignments matters too: primates4.chr6 aligns
about as many bases per rank as yeast7 but in 4.5$\times$ as many
alignments, and takes 1.7$\times$ as long.

Second, the stage lasts as long as its slowest rank, and the work is
unevenly assigned (see above). On
primates14.chr6 the slowest rank takes 3.2$\times$ the average, which
together with the first factor accounts for its 28$\times$ growth.
